\chapter*{Заключение}
\addcontentsline{toc}{chapter}{Заключение}

В рамках работы была рассмотрена задача генерации коллективных музыкальных рекомендаций для посетителей заведений и мероприятий методами машинного обучения. Работа сочетает бизнес- и продуктовую проработку сервиса, анализ рынка и правовой среды, а также техническую часть, связанную с моделированием групповых рекомендаций на датасете YAMBDA.

В первой главе исследована предметная область музыкального сопровождения заведений и мероприятий. Показано, что музыка в коммерческих и событийных пространствах выполняет не только фоновую, но и сервисную, маркетинговую и сценарную функцию. Были проанализированы сегменты рынка, существующие практики музыкального сопровождения, конкурентные B2B-сервисы и правовые аспекты публичного воспроизведения музыки. Отдельно рассмотрены требования к взаимодействию с РАО и ВОИС и необходимость отчетности по воспроизведенным произведениям и фонограммам. Рыночная оценка показала, что постоянные B2B-точки и разовые события следует считать разными типами рынка; базовый адресуемый рынок для постоянных точек оценен примерно в 25 тыс. объектов, уже находящихся в легальном аудио-контуре, а событийный сценарий~--- примерно в 12 тыс. лицензируемых мероприятий в год.

Во второй главе разработана концепция сервиса «Резонанс» / «Яндекс Музыка Резонанс». Сервис описан как часть экосистемы «Яндекс Музыки», в которой бизнес получает специальный профиль, задает рамки музыкального потока, а посетители добровольно участвуют в формировании групповой «волны». В концепции предусмотрены BLE-beacon как основной механизм фиксации присутствия, QR-код как дополнительный сценарий, обработка данных внутри контура Яндекса, информированное согласие пользователя, отсутствие передачи персональных данных заведению, настройка музыкальных ограничений, возможный AI-консультант и автоматическая генерация отчетности для РАО и ВОИС. Также рассмотрены дополнительные применения технологии, включая помощника DJ и анализ вовлеченности аудитории по зонам пространства.

В третьей главе систематизированы теоретические и технологические основы построения музыкальных рекомендательных систем. Рассмотрены классическая коллаборативная фильтрация, последовательные модели, трансформерные архитектуры семейства SASRec/BERT4Rec/gSASRec, специфика музыкального домена, повторное потребление и использование аудиоэмбеддингов. Эта часть обосновывает выбор персонального скорера и особенности протокола оценки, применяемого в экспериментальной части.

В четвертой главе рассмотрены подходы к формированию коллективных музыкальных рекомендаций. Были описаны классические стратегии агрегации персональных оценок, различие между persistent- и эфемерными группами, а также нейросетевые методы групповых рекомендаций, включая AGREE, GroupIM и AlignGroup. На этой основе сформирован ландшафт бейзлайнов и архитектурных вариантов, необходимых для экспериментального сравнения.

В пятой главе реализована и исследована модель коллективных музыкальных рекомендаций на основе датасета YAMBDA. Экспериментальная часть демонстрирует, как из индивидуальных историй и персонального скорера формируются кандидаты для группы, как сравниваются классические и обучаемые агрегаторы и как оценивается качество групповых рекомендаций. Особое внимание уделено использованию аудиоэмбеддингов как индивидуального контентного сигнала, который может быть поднят на групповой уровень.

В шестой главе экспериментальная модель была связана с демонстрационным сценарием. На уровне алгоритмического прототипа показано, как список идентификаторов пользователей из YAMBDA может задавать текущую эфемерную группу, для которой пересчитывается ранжированный список анонимизированных треков и метрики качества. Отдельно рассмотрен Android BLE-прототип, демонстрирующий локальное определение ближайшего сохраненного маяка по относительной силе Bluetooth-сигнала. Этот компонент не является системой точной геолокации, но иллюстрирует техническую возможность использовать BLE-сигнал как приближенный индикатор близости к зоне пространства.

Работа имеет несколько существенных ограничений. Во-первых, авторы не располагают доступом к реальным внутренним данным «Яндекс Музыки», промышленным алгоритмам рекомендаций и B2B-лицензионной инфраструктуре платформы. Во-вторых, датасет YAMBDA является обезличенным, поэтому невозможно анализировать реальные названия рекомендованных треков, жанровые интерпретации конкретных списков и пользовательские профили в человекочитаемом виде. В-третьих, промышленная реализация сервиса описана концептуально: в рамках ВКР не реализуются получение пользовательских согласий, интеграция с реальным бизнес-профилем, юридически полноценный контур публичного воспроизведения и промышленная BLE-инфраструктура. Android-приложение рассматривается только как демонстрация принципа работы с сохраненными BLE-маяками и RSSI-сигналом. В-четвертых, рыночные оценки для событийного сегмента являются сценарными, поскольку открытая статистика не содержит достаточно узкой строки по вечеринкам, DJ-сетам и специальным мероприятиям.

Перспективы дальнейшего развития связаны с несколькими направлениями. На продуктовой стороне требуется проведение дополнительных интервью и пилотов с владельцами заведений, организаторами мероприятий, DJ и пользователями, чтобы проверить готовность к использованию сервиса и уточнить ценностное предложение. На правовой стороне необходима детальная проработка B2B-лицензирования музыкального каталога, форматов отчетности и пользовательского согласия. На технической стороне перспективны фильтрация и бизнес-ограничения при формировании плейлиста, неравномерное взвешивание вклада пользователей, использование нескольких BLE-зон, устойчивость к холодному старту, fairness-метрики для групп и онлайн-обновление рекомендаций. Отдельными направлениями могут стать разработка ассистента для DJ и формирование B2B-аналитики на основе новых данных о музыкальном поведении аудитории. Эти направления следует рассматривать как гипотезы дальнейших исследований и продуктовой проверки, а не как доказанный бизнес-эффект текущего прототипа.

Таким образом, в работе показано, что сервис коллективных музыкальных рекомендаций для заведений и мероприятий может быть рассмотрен как содержательная продуктовая ниша между классическими B2B-сервисами фоновой музыки и сервисами индивидуального заказа треков. Его ценность определяется сочетанием трех элементов: учета фактического состава аудитории, сохранения управляемости музыкальной атмосферы для бизнеса и соблюдения правовых и приватностных ограничений. Экспериментальная часть на YAMBDA, в свою очередь, демонстрирует техническую основу для дальнейшего исследования групповых музыкальных рекомендаций в таком сценарии.
